\chapter{Failure to Thrive}

\section*{Definition}
Failure to thrive (FTT) is a state of inadequate physical growth and development, typically defined as weight persistently below the 3rd percentile, weight deceleration crossing two major percentile lines, or weight less than 80\% of ideal body weight for age. It reflects a complex interplay of medical, nutritional, and psychosocial factors.

\section*{What Happens in Your Body}
FTT results from insufficient caloric intake, excessive energy expenditure, impaired nutrient absorption, or a combination thereof. Inadequate nutrition disrupts anabolic pathways, leading to reduced insulin-like growth factor 1 (IGF-1) signaling, catabolism of lean body mass, and stunting of linear growth. Chronic undernutrition impairs immune function, neurodevelopment, and organ system maturation, particularly in the first 2--3 years of life.

\section*{Causes}
\begin{itemize}
  \item \textbf{Inadequate intake:} Poor feeding technique, breastfeeding difficulties, oral-motor dysfunction, poverty, food insecurity, neglect
  \item \textbf{Increased losses:} Gastroesophageal reflux, vomiting, diarrhea, malabsorption syndromes (celiac disease, cystic fibrosis)
  \item \textbf{Increased metabolic demand:} Congenital heart disease, chronic infection, hyperthyroidism, malignancy
  \item \textbf{Organic disease:} Cerebral palsy, chromosomal abnormalities, lead poisoning, renal tubular acidosis, hypothyroidism
\end{itemize}

\section*{When to See a Doctor}
See your doctor if:
\begin{itemize}
  \item Weight loss or flat growth trajectory over 2--3 consecutive well-child visits
  \item Accompanying developmental regression, vomiting, diarrhea, or feeding refusal
  \item Signs of dehydration, pallor, or wasting
  \item Concerns regarding parental mental health, substance use, or social environment
\end{itemize}

\section*{Related Symptoms}
\begin{itemize}
  \item Short stature, developmental delay, irritability, lethargy, muscle wasting, recurrent infections
\end{itemize}
\textbf{Important:} A multidisciplinary approach involving pediatrics, nutrition, social work, and occupational therapy is essential to address both medical and environmental contributors to FTT.

\section*{Self-Care / Home Management}
\begin{itemize}
  \item Keep a 3--7 day food diary to assess caloric intake and feeding patterns
  \item Offer nutrient-dense, age-appropriate foods at structured meal and snack times
  \item Minimize distractions during feeding; avoid juice and sugar-sweetened beverages
\end{itemize}

\section*{How Your Doctor Will Evaluate You}
\begin{itemize}
  \item Plot growth parameters on standardized WHO or CDC growth charts; calculate weight-for-length or BMI
  \item Detailed dietary, perinatal, developmental, and psychosocial history
  \item Directed laboratory evaluation: CBC, CMP, ESR, celiac panel, thyroid function, lead level, sweat chloride test
\end{itemize}

\section*{Treatment Options}
\begin{itemize}
  \item Nutritional intervention: Caloric supplementation, high-energy formulas, feeding therapy with occupational or speech therapist
  \item Treat underlying medical conditions such as GERD, celiac disease, or infection
  \item Social support: Referral to WIC, SNAP, early intervention, and social work for environmental barriers
  \item Regular follow-up with weight checks every 1--4 weeks until growth trajectory normalizes
\end{itemize}

\section*{References}
\begin{thebibliography}{99}
\bibitem{ref1} Kliegman RM, St Geme JW III, Blum NJ, et al. "Nelson Textbook of Pediatrics." 22nd ed. Elsevier; 2024.
\bibitem{ref2} Cole SZ, Lanham JS. "Failure to Thrive: An Update." Am Fam Physician. 2011;83(7):829--834.
\bibitem{ref3} Gahagan S. "Failure to Thrive: A Consequence of Undernutrition." Pediatr Rev. 2006;27(1):e1--e11.
\bibitem{ref4} Perrin EC, Cole CH, Frank DA, et al. "Criteria for Determining Disability in Infants and Children: Failure to Thrive." Pediatrics. 2003;111(1):177--184.
\bibitem{ref5} Wright CM. "Identification and Management of Failure to Thrive." Arch Dis Child. 2000;83(1):78--81.
\end{thebibliography}
\clearpage
